Failure to do so risks conflating pharmacological artifacts with biological time.  

The distinction is fundamental. Biological time is the calendar interval during which a viral lineage replicates under the endogenous error rate of its RNA-dependent RNA polymerase, subject to natural selection, genetic drift and recombination. Pharmacological artifacts are the excess substitutions—predominantly G-to-A and C-to-U transitions—generated by NHC-TP incorporation during molnupiravir treatment. These substitutions accumulate in days rather than years, carry a diagnostic spectral signature, and can be transmitted into descendant clades. When phylogenetic methods treat every observed nucleotide difference as a product of biological time, the artifacts are silently converted into inflated estimates of evolutionary rate, erroneous divergence dates and spurious inferences about silent circulation or multi-host cascades. The thesis holds that this conflation is not a minor statistical inconvenience; it is a systematic distortion that undermines the temporal calibration of the SARS-CoV-2 phylogeny after 2021 and, by extension, every historical claim that depends on that calibration.

Consider first the mechanics of the artifact. In an untreated host the substitution rate remains on the order of \(10^{-3}\) to \(10^{-4}\) per site per year. Under NHC-TP pressure the rate rises approximately eight-fold and becomes heavily biased toward transitions. A single treated infection can therefore generate a long phylogenetic branch containing dozens to more than a hundred excess mutations within a clinical window of days to weeks. When that genome is transmitted, the excess distance is inherited by every descendant. In a global tree these events appear as abrupt rate accelerations or as long terminal branches whose lengths, if interpreted under a constant or slowly relaxing clock, imply either extended periods of unobserved circulation or sustained elevation of the endogenous mutation rate. Neither interpretation is correct. The distance is real, yet the time required to produce it is pharmacological, not biological. Failure to flag the spectral signature—high G-to-A / C-to-T ratio, preferred nucleotide contexts, geographic and demographic correlation with prescription data—allows the artifact to be absorbed into the clock model as ordinary evolutionary change.

The consequences propagate in both temporal directions. Backward extrapolation across the 2022 threshold attempts to recover the timing of early human emergence or to test the feasibility of a multi-species cascade inside the narrow window that followed the end of intensive sampling on 31 May 2019. If a fraction of the post-2022 genetic distance is pharmacological, any clock calibrated on the contaminated data will over-estimate the rate that operated in 2019–2021. The inferred tMRCA is pushed earlier, the apparent duration of pre-detection circulation is lengthened, and the multiplicative improbability of an unobserved natural proximal ancestor carrying the exact PRRA + CGG-CGG configuration is correspondingly relaxed. Conversely, analyses that ignore the artifact and treat the entire post-2022 tree as clock-like under-estimate the mutational supply available for adaptation once signature-bearing lineages are already circulating. Forward projections of diversity, of the probability of further antigenic change, or of the persistence of the exhaustion attractor become systematically optimistic.

The same conflation distorts tests of irreversibility. The original argument rested on the continuous advance of an endogenous molecular clock: any 2019 lineage that continued to replicate would, by 2026, differ from its progenitor by six to seven years of ordinary substitutions. The pharmacological layer adds a second, non-uniform increment. When that increment is misread as biological time, the genetic distance separating modern specimens from a hypothetical 2019 proximal ancestor is still greater than the pure-clock prediction, yet the excess is attributed to natural processes rather than to intermittent mutagenic pressure. The practical impossibility of recovering a temporally authentic 2019 genome is thereby masked; the zero probability that follows from irreversible evolutionary time is diluted by an artifactual explanation that appears to restore a longer, more permissive natural history.

Within the computational formulation the risk is rendered quantitative. Year-specific matrices encode the progressive elevation of the mutagenic rate, the dominance of the directional spectrum, and the coupled decline in codon-usage bias together with rises in titer, glycosylation modification and brain tropism. Each simulated generation adds mutations whose composition matches the pharmacological signature. Because the model advances through successive years without resetting the spectral bias, the cumulative distance between the 2019 baseline and the 2026 state systematically exceeds the distance expected under a pure endogenous clock. An analyst who fitted a standard molecular-clock model to the simulated sequences without accounting for the NHC-TP term would recover an erroneously accelerated rate, an inflated root-to-tip slope, and a set of node ages that no longer correspond to calendar time. The numerical experiment therefore demonstrates the precise mechanism of conflation: pharmacological increments are absorbed into the biological-time parameter, and every downstream inference is shifted.

The practical requirement that follows is methodological, not ideological. Post-2022 genomes must be screened for the molnupiravir spectral signature before they are admitted into rate estimation, dating or comparative analyses that span the 2022 threshold. Branches that meet the high G-to-A / C-to-T criteria, that cluster geographically with known prescription intensity, or that seed small descendant clades carrying the same imprint should be flagged, down-weighted or modelled with an explicit pharmacological rate component. Failure to apply that filter does not merely introduce noise; it systematically substitutes a chemically induced mutational process for the biological time that phylogenetic reconstruction is intended to recover. The thesis therefore concludes that the integrity of temporal inference after 2021 depends on the explicit separation of pharmacological artifacts from endogenous evolutionary change. Without that separation the clock is not merely broken; its fragments are reassembled into a narrative that no longer measures the history it claims to describe.
